% !TEX root = ../main.tex

\chapter{相关技术与理论基础}

\section{技术综述的聚焦范围}

围绕长期运行的 Flink 作业在 Sink 侧动态组织 Kafka 写入路径这一目标，后续需求分析与系统设计主要依赖三类基础：
Kafka 写入端的 Topic 元数据、事务与可见性语义，Flink 的 Checkpoint 与统一 Sink API，
以及现有开源连接器和社区提案在动态目标管理方面的能力边界。
因此，本章从这些内容展开，为后文分析动态 Sink 需求、设计 route 级 writer 管理和提交流程提供技术依据。

\section{Kafka 写入端相关机制}

\subsection{Topic 元数据与管理面访问}

Apache Kafka 的 Topic 是生产端写入目标的基本单位\cite{kreps2011kafka}。对于本文所研究的动态 Sink 而言，
关键并不是 Kafka 的完整消息模型，而是两个与后续设计直接相关的事实：第一，Topic 名称能够被用作动态写入目标的物理标识；
第二，Topic 集合可以通过管理面接口在作业运行期间被重新查询。Kafka Connect、Debezium 等开源项目也依赖 Kafka
作为数据交换底座，但它们通常把 Topic 当作 connector 配置或外部系统映射的结果\cite{kafkaConnectDocs,debeziumDocs}；
本文则需要把 Topic 集合作为运行期可变化的元数据输入，驱动 Sink writer 的创建、复用和回收。
当前原型正是利用
\id{AdminClient.listTopics()} 获取集群中可见的 Topic 名称，再把这些名称封装为
\id{KafkaStream}，供后续正则匹配与 route 解析使用。

因此，Kafka 在本文中同时扮演两种角色：一方面，它是承接 Flink 输出结果的下游消息系统；
另一方面，它又是一个需要被动态感知的外部元数据源。第 4 章中
\id{SingleClusterTopicMetadataService} 的职责，正是把 Kafka 管理面暴露出的 Topic 列表转换为连接器内部统一使用的
\id{KafkaStream} / \id{ClusterMetadata} 结构。这一点决定了本文的“动态发现”发生在 Topic 粒度，
而不是 partition 粒度或消费者组订阅粒度。

\begin{figure}[htbp]
  \centering
  \safeincludegraphics[width=0.9\textwidth]{kafka_flink_pipeline.png}
  \caption{Kafka 与 Flink 组成的典型实时数据流水线}
  \label{fig:kafka-flink-pipeline}
\end{figure}

图~\ref{fig:kafka-flink-pipeline} 展示了 Kafka 与 Flink 在实时数据链路中的典型位置：
多类数据源先汇入 Kafka，Flink 再作为中间处理与分析引擎，将结果写入下游数据汇聚系统或业务系统。
本文所研究的动态 Sink，正位于这一链路的“Flink 写出下游”位置，其目标不是替代 Kafka 或 Flink，
而是在二者集成边界上增强写入端的动态适配能力。

\subsection{Topic、Partition 与写入目标绑定}

Topic 是逻辑上的消息类别，其下可划分为多个分区以实现并行与扩展。
消息在分区内有序，分区之间相互独立。动态 Sink 的设计重点并不是重新实现 Kafka 的分区调度，
而是决定每条记录最终交给哪个 Topic 以及哪个底层 producer/writer 实例。因此，本文把 Topic 作为动态 route 的物理落点：
自动发现路径将 \id{stream_pattern} 命中的 Topic 展开为 \id{activeRouteIds}；
显式路由路径则把业务 \id{routeId} 解析为一个或多个实际 Topic。

一旦 route 被解析为具体 Topic，底层 \id{KafkaSink} 仍会按照 Kafka producer 的常规逻辑进行序列化、
分区选择和发送确认。换言之，本文并不改变 Kafka 对 partition 的写入语义，而是在 partition 选择之前增加了一层
“逻辑 route $\rightarrow$ 物理 Topic”的动态绑定。这个分层关系解释了为什么第 4 章需要在
\id{DynamicKafkaSinkWriter.createClusterWriter(...)} 中为每个物理 route 构造独立的底层 writer：
只有先把动态 route 解析成确定 Topic，后续 Kafka producer 才能按原有机制完成真正发送。

这一边界也说明了本文与 Kafka 分区优化类工作的差异。分区负载均衡或性能建模研究关注的是一个 Topic 内部或集群节点之间的负载分布
\cite{wang2022astraea,wu2021logp}，而本文关注的是写入目标集合本身会变化时，Flink Sink 如何在 Topic 粒度上重新组织 writer。
换言之，本文并不试图替代 Kafka 的分区器，而是在 Kafka 分区器发挥作用之前，先解决“本条数据应该进入哪些 Topic”。

\subsection{生产者事务、精准一次语义与可提交读}

本文后续涉及 \id{EXACTLY_ONCE} 验证，因此需要明确 Kafka 写入端一致性语义的边界。Kafka 的普通 producer
只能保证记录成功发送到 broker，并不能天然避免故障重试导致的重复写入。为减少重试重复，Kafka 引入了幂等 producer；
而要把一组写入作为事务提交或回滚，则需要配置 \id{transactional.id}，由 producer 在事务边界内写入记录，
并最终执行 commit 或 abort\cite{kafkaSemantics}。

事务提交后，Kafka 消费端还需要通过隔离级别决定是否读取未提交数据。\id{isolation.level=read_committed}
表示消费者只读取已经提交的事务数据以及非事务数据；\id{read_uncommitted} 则可能读到尚未提交、后续也可能被回滚的记录。
这一点对本文实验具有直接意义：如果只查看 topic offset 或使用默认消费配置，可能只能说明 broker 收到了数据，
但不足以严格证明事务数据已对下游以 committed 形式可见。因此，第 5 章在说明事务路径可运行性验证时，
关注的是 Checkpoint 成功、事务提交未报错以及目标 Topic 出现 committed 输出，而不是把其等同于完整故障注入下的端到端精准一次证明。

对于动态 Sink 来说，Kafka 事务还带来额外复杂性：静态 Sink 通常只需要围绕固定 Topic 和固定事务前缀维护 producer；
动态 Sink 则可能在运行期新增或移除 route。如果多个物理 route 共用同一事务标识，提交和恢复时容易混淆不同目标的事务边界。
因此，本文实现采用 route 级事务前缀派生策略，即由全局 \id{transactionalIdPrefix} 和 \id{routeId} 生成不同物理 route 的事务前缀。
第 4 章中 \id{buildTransactionalIdPrefix(routeId)}、\id{DynamicKafkaCommittable} 与
\id{DynamicKafkaCommitter} 的设计，正是建立在这一 Kafka 事务模型之上。

从现有能力看，Kafka 已经提供事务 producer 与 \id{read_committed} 消费隔离级别，
Flink KafkaSink 也已经能够在静态目标场景下利用这些能力实现事务提交\cite{flinkKafkaDocs,kafkaSemantics}。
但这些能力默认并不回答“运行期新增 Topic 对应的 producer 何时创建、事务前缀如何命名、旧 route 的事务状态何时能够安全释放”等问题。
这正是动态 Sink 相对静态 Sink 的特殊需求。

\section{Flink 容错与 Sink 提交机制}

\subsection{Checkpoint 与动态连接器状态}

Flink 通过 Checkpoint 与分布式快照思想实现容错\cite{carbone2015flink}。
在本文场景中，Checkpoint 需要解决的不是普通算子状态如何保存这一通用问题，而是动态连接器运行期产生的 route 状态能否被恢复。
动态 Sink 会在运行期间维护 \id{activeRouteIds}、\id{clusterWritersByRouteId}、显式 route 表以及底层 Kafka writer 状态。
如果这些信息不进入快照，作业重启后就可能丢失已经创建的写入通道，或者无法继续提交故障前产生的事务。

因此，当前代码将动态状态组织为 route 级恢复单元：\id{DynamicKafkaSinkWriterState} 记录 route 标识、
目标 cluster、Topic、事务前缀、producer 配置和底层 \id{KafkaWriterState}；
\id{DynamicKafkaSink.restoreWriter(...)} 再利用这些状态重建 route 对应的底层 writer。
这说明后文的状态设计不是为了形式上“支持 Checkpoint”，而是为了让动态 route 在故障恢复后仍能与 Kafka 写入目标重新对齐。

\subsection{统一 Sink API 与两阶段提交}

Flink 统一 Sink API 将写入逻辑拆分为 Writer、状态快照、提交前准备和 Committer 等阶段\cite{flinkKafkaDocs,flinkUnifiedSink}。
在 \id{AT_LEAST_ONCE} 语义下，Writer 通常只需保证数据被刷出到外部系统；在 \id{EXACTLY_ONCE} 语义下，
Sink 还需要与 Checkpoint 配合，把某个 Checkpoint 之前的写入作为可提交单元交给 Committer。
这也是本文没有简单地在用户函数里手动创建 Kafka producer 的原因：手写 producer 很难自然接入 Flink 的状态快照和提交回调。

当前实现中，\id{DynamicKafkaSink} 直接实现
\id{TwoPhaseCommittingStatefulSink<IN, DynamicKafkaSinkWriterState, DynamicKafkaCommittable>}。
这一选择使动态 Sink 能够在 Writer 阶段解析 route 并写入数据，在 \id{snapshotState(...)} 中保存 route 级 writer 状态，
在 \id{prepareCommit()} 中把底层 \id{KafkaCommittable} 包装为带 route 信息的
\id{DynamicKafkaCommittable}，最后由 \id{DynamicKafkaCommitter} 按 route 分组提交。
第 4 章对状态和提交路径的设计，正是对这一 API 分阶段语义的具体展开。

需要强调的是，统一 Sink API 提供的是扩展连接器的框架，而不是动态 Kafka Sink 的完整实现。
官方 KafkaSink 已经很好地覆盖了“固定目标 Topic + 指定投递语义”的主流场景\cite{flinkKafkaDocs}，
但它并不会在作业运行期间根据 Topic 正则自动扩展目标集合，也不会为逻辑 route 维护多组底层 writer。
因此，本文的研究独特性不在于重新发明 Sink API，而在于把该 API 的 Writer、状态和 Committer 阶段重新组织到 route 粒度，
使动态目标管理能够纳入 Flink 原有的容错和提交框架。

\subsection{控制流、数据流与广播状态}

本文的动态 route 验证依赖应用层广播状态：\id{DynamicJob} 读取 YAML 中的 \id{route_map}，
生成 \id{DynamicSinkEvent.update(...)} 控制事件，再通过 \id{BroadcastProcessFunction} 将 route 配置写入广播状态并转发给 Sink。
这一机制本身属于 Flink DataStream 层的能力，但它与 connector 设计密切相关：广播状态负责在进入 Sink 前过滤无效
\id{routeId}，Sink 内部则负责解释 route 更新事件并维护 \id{explicitRoutesById}。

因此，本文中的控制流与数据流不是两个完全独立的系统。控制事件和普通数据事件最终都会进入
\id{DynamicKafkaSinkWriter.write(...)}，再由 Writer 根据 \id{DynamicKafkaRouteEvent} 的语义分流。
这一点为第 3 章“控制面配置必须先于数据面生效”的需求分析提供了基础，也解释了第 4 章为什么要把
\id{DynamicKafkaRouteEvent} 设计为 Sink 可识别的统一事件接口。

\section{Flink Kafka Connector}

\subsection{静态 Kafka Sink 的目标绑定}

常规配置在作业构建阶段指定固定的 Topic 名称或静态列表，运行期不随集群侧 Topic 集合变化而自动扩展。
该方式实现简单、语义清晰，适用于拓扑稳定的场景，但在 Topic 频繁增删或按规则批量创建时，
难以避免停启作业与重复发布。本文原型中的对照组也可以理解为这种典型静态模式：
应用在作业提交前就把目标 Topic 写死到 \id{KafkaSink} 中，之后任何目标变更都需要重新构建并部署作业。
因此，静态模式更适合作为本文后续实验中的基准方案，而不是待研究问题的最终答案。

从能力边界看，静态 KafkaSink 的优势是成熟、稳定、语义清楚，尤其适合“目标已知且长期稳定”的作业。
但本文的实际需求恰好相反：目标 Topic 可能随业务上线、租户扩展或路由规则变化而出现，
且用户希望 Flink 作业不因目标变化而停机。静态 KafkaSink 无法在不重启作业的前提下重新扫描 Topic 集合，
也没有内建逻辑 route 到物理 Topic 的映射表。因此，本文把它作为复用对象和对照基线，而不是直接采用其原始形态。

\subsection{正则发现从 Source 侧到 Sink 侧的差异}

在 Source 侧，社区已支持基于正则表达式匹配 Topic 名称的订阅方式，并与分区发现、
偏移提交等机制结合\cite{flip246}。其本质是周期或事件驱动地比对“模式—当前元数据”，
更新实际订阅集合。Sink 侧若要对齐类似能力，需要另行设计生产者实例生命周期与记录路由策略，
不能简单照搬 Source 实现。本文当前代码中对“按模式发现 Topic”的实现落在
\id{DynamicKafkaSinkBuilder.setStreamPattern(...)}、\id{StreamPatternSubscriber} 与
\id{SingleClusterTopicMetadataService} 的组合上；但当模式命中新的 Topic 后，
系统还要进一步将其转换为 route id、Kafka producer 配置和底层 writer，
这一步正是 Sink 侧相较 Source 侧的主要额外复杂度。

为了更直观地对应本仓库中的实现逻辑，代码清单~\ref{code:metadata-and-pattern}
展示了当前原型中最关键的两步：先通过 \id{AdminClient.listTopics()} 拉取所有 Topic，
再通过正则对 \id{streamId} 执行匹配。这说明本文方案中的“动态发现”并不是来自
Flink 运行时的内建事件通知，而是构建在 Kafka 管理面元数据轮询之上的。

\begin{codeblock}[language=Java,caption={当前原型中的 Topic 元数据获取与正则匹配逻辑},label={code:metadata-and-pattern}]
// connector/.../SingleClusterTopicMetadataService.java
public Set<KafkaStream> getAllStreams() {
    return getAdminClient().listTopics().names().get().stream()
            .map(this::createKafkaStream)
            .collect(Collectors.toSet());
}

// connector/.../StreamPatternSubscriber.java
public Set<KafkaStream> getSubscribedStreams(KafkaMetadataService kafkaMetadataService) {
    Set<KafkaStream> allStreams = kafkaMetadataService.getAllStreams();
    List<KafkaStream> matches = new ArrayList<>();
    for (KafkaStream kafkaStream : allStreams) {
        String streamId = kafkaStream.getStreamId();
        if (streamPattern.matcher(streamId).find()) {
            matches.add(kafkaStream);
        }
    }
    return Set.copyOf(matches);
}
\end{codeblock}

结合代码可以看到，当前实现默认把 \id{streamId} 等价为 Topic 名称，因此
\id{stream_pattern} 实际上就是“对 Topic 名称做匹配”的正则表达式。这样做实现简单，
也便于原型阶段验证功能；但如果未来扩展到“逻辑流名称”与“物理 Topic 名称”并不完全一致的场景，
则需要在 \id{KafkaStream} 的建模上进一步引入更丰富的元数据层次。

\subsection{开源数据接入项目的能力边界}

除 Flink Kafka Connector 外，Kafka 生态中还存在 Kafka Connect、Debezium、Flink CDC 等成熟项目。
Kafka Connect 提供 source connector 与 sink connector 框架，用于在 Kafka 与外部系统之间搬运数据\cite{kafkaConnectDocs}；
Debezium 通过读取数据库变更日志，将表级变更持续写入 Kafka\cite{debeziumDocs}；
Flink CDC 则把变更数据捕获能力接入 Flink 作业，用于实时同步和实时数仓构建\cite{flinkCdcDocs}。
这些项目说明，开源生态已经能够很好地解决“如何把外部系统的数据接入流式链路”这一类问题。

但它们与本文需求存在本质区别。Kafka Connect 和 Debezium 的动态性主要体现在 connector 任务配置、
外部表捕获范围或任务实例管理上；Flink CDC 关注的是数据库变更如何成为 Flink 输入流。
本文则研究 Flink 作业已经在运行、数据已经进入计算链路之后，Sink 侧如何根据 Kafka Topic 元数据变化动态改变写出目标。
这一问题要求 writer 生命周期、route 状态、Kafka 事务与 Flink Checkpoint 同时协同，
并不是简单增加一个外部 connector 配置项即可解决。

\begin{table}[htbp]
  \centering
  \caption{现有方案能力边界与本文需求对比}
  \label{tab:related-capability-gap}
  \small
  \begin{tabular}{p{0.17\textwidth}p{0.29\textwidth}p{0.36\textwidth}}
    \toprule
    方案或工作 & 已具备能力 & 与本文需求的差异 \\
    \midrule
    KafkaSink & 固定 Topic 写入、投递语义、事务提交 & 目标 Topic 通常在构建阶段绑定，缺少运行期正则发现和 route 映射 \\
    FLIP-246 & Source 侧动态发现 Topic 与分区订阅 & 面向消费端 offset 和分区管理，不能直接解决 producer、writer 和提交问题 \\
    FLIP-515 & 提出动态 Sink 方向与接口讨论 & 公开资料以设计讨论为主，仍需要可运行原型和实验验证支撑 \\
    Kafka Connect 等 & 外部系统接入、CDC 与数据同步 & 关注数据进入 Kafka 或 Flink，不负责 Flink Sink 内部的动态多目标写出 \\
    \bottomrule
  \end{tabular}
\end{table}

\subsection{动态 Sink 需要解决的连接器问题}

现有 Kafka Sink 在统一 API 下仍以静态目标为主；FLIP-515\cite{flip515} 等工作讨论动态 Sink 的方向，
但生产可用的完整方案需处理：元数据扫描开销、写入路径与事务语义、状态快照与恢复后的对账等问题。
从当前原型实现看，这些难点都已经在代码层面有所体现：例如 route 级 writer 的惰性创建已经实现，
非事务路径下的退役 writer 清理也已经落地；但 \id{EXACTLY_ONCE} 路径下的安全回收、
事务边界与复杂故障验证仍待补足。
因此，Flink Kafka Connector 的“动态化”并不是单点功能，而是一组围绕 Writer、状态、提交与元数据服务的协同扩展。

结合表~\ref{tab:related-capability-gap} 可以看到，现有方案已经分别解决了静态 Kafka 写入、Source 侧动态订阅、
外部系统接入和动态 Sink 方向讨论等问题，但没有同时覆盖本文所需的四个约束：
第一，运行期根据 Kafka Topic 元数据发现目标；第二，支持业务逻辑 route 与物理 Topic 解耦；
第三，以 route 为粒度维护底层 writer、状态和提交对象；第四，在原型阶段通过实验验证发现实时性、吞吐、Checkpoint 与事务路径。
本文的独特性正来自这四个约束的组合，而不是来自某一个单独技术点。

\section{本章小结}

本章围绕后续设计所需的关键技术进行了有针对性的综述。Kafka 部分重点讨论 Topic 元数据、
写入目标绑定、事务型 producer、Exactly-once 与 \id{read_committed} 可见性；
Flink 部分重点说明 Checkpoint、统一 Sink API、两阶段提交和广播状态如何支撑动态连接器；
Connector 部分则对比了静态 KafkaSink、FLIP-246、FLIP-515、Kafka Connect、Debezium 与 Flink CDC 等已有工作。
可以看到，现有方案已经提供了若干可复用能力，但无法完整满足本文“Sink 侧运行期动态发现、多 route 写出、状态快照和提交协同”的组合需求。
因此，本文后续提出的动态 Kafka Sink 原型并非对通用技术的重复介绍，
而是在现有能力边界之上面向特定场景进行的连接器扩展。
